% Figure: wave-equation snapshot showing a triangular pulse on the
% line decomposing into a left-moving and right-moving half-amplitude
% copy, per d'Alembert. Three snapshots in time.
\begin{figure}[htbp]
\centering
\begin{tikzpicture}
\begin{axis}[
  width=12cm, height=6cm,
  xlabel={position $x$},
  ylabel={displacement $u(x,t)$},
  xmin=-5, xmax=5,
  ymin=-0.1, ymax=1.1,
  grid=both,
  major grid style={engblack!20},
  minor grid style={engblack!10},
  legend pos=north east,
  legend cell align=left,
  font=\small,
  samples=300,
  domain=-5:5,
]
% Initial triangular pulse centred at origin, half-width 1.
\addplot[engblack, very thick, dashed] {max(0, 1 - abs(x))};
\addlegendentry{$t = 0$: triangular pulse}

% After short time: two half-amplitude copies at +/- 1.5
\addplot[engblue, very thick] {
   0.5*max(0, 1 - abs(x-1.5))
 + 0.5*max(0, 1 - abs(x+1.5))
};
\addlegendentry{$t = 1.5/c$: split}

% After longer time: half-amplitude copies at +/- 3
\addplot[engred, very thick] {
   0.5*max(0, 1 - abs(x-3))
 + 0.5*max(0, 1 - abs(x+3))
};
\addlegendentry{$t = 3/c$: separated}

% Characteristics: dotted lines x = +/- c t through the peak
\draw[engblue!50, dotted, thick] (axis cs:0,0) -- (axis cs:3,0.5);
\draw[engblue!50, dotted, thick] (axis cs:0,0) -- (axis cs:-3,0.5);

\end{axis}
\end{tikzpicture}
\caption[Wave-equation d'Alembert decomposition]{Wave-equation
solution on the infinite line for a triangular initial pulse
released from rest. By the d'Alembert formula
$u(x,t) = \tfrac{1}{2}(\phi(x-ct) + \phi(x+ct))$, the pulse splits
into a half-amplitude right-moving copy and a half-amplitude
left-moving copy, each translating at speed $c$ without changing
shape. The dotted blue lines trace the characteristics $x = \pm c
t$ along which information propagates. Finite propagation speed is
the working signature of hyperbolic equations and the formal
content of the bounded domain of dependence: $u(x,t)$ depends only
on initial data in $[x-ct, x+ct]$.}
\label{fig:vol02:ch12:wave-snapshot}
\end{figure}
